% P10-4c, 14.09.2026: Ledger-Herleitung und Nachweisanschlüsse präzisiert.
% B-46-Nachtrag, 13.09.2026: Z3 an die Originalprüfung E43-05
% und die bereits synchronisierten Fassungen von 4.3/Anhang A angeglichen.
% M16 / B-34, 11.09.2026: Z4 trennt die Artefakte der beiden Prüfbeobachtungen.
% M03 / B-17, 10.09.2026: ursprüngliches und erweitertes Nutzungsziel
% unterschieden; keine neue Datierung des Betreuerimpulses.
% ============================================================
% Tabelle zu §4.2 — Ziel gemäß aktueller Einbindung in chap4.2.tex:
%   tables/chap4/chap4.2table.tex
% Aufruf: % P10-4c, 14.09.2026: Ledger-Herleitung und Nachweisanschlüsse präzisiert.
% B-46-Nachtrag, 13.09.2026: Z3 an die Originalprüfung E43-05
% und die bereits synchronisierten Fassungen von 4.3/Anhang A angeglichen.
% M16 / B-34, 11.09.2026: Z4 trennt die Artefakte der beiden Prüfbeobachtungen.
% M03 / B-17, 10.09.2026: ursprüngliches und erweitertes Nutzungsziel
% unterschieden; keine neue Datierung des Betreuerimpulses.
% ============================================================
% Tabelle zu §4.2 — Ziel gemäß aktueller Einbindung in chap4.2.tex:
%   tables/chap4/chap4.2table.tex
% Aufruf: % P10-4c, 14.09.2026: Ledger-Herleitung und Nachweisanschlüsse präzisiert.
% B-46-Nachtrag, 13.09.2026: Z3 an die Originalprüfung E43-05
% und die bereits synchronisierten Fassungen von 4.3/Anhang A angeglichen.
% M16 / B-34, 11.09.2026: Z4 trennt die Artefakte der beiden Prüfbeobachtungen.
% M03 / B-17, 10.09.2026: ursprüngliches und erweitertes Nutzungsziel
% unterschieden; keine neue Datierung des Betreuerimpulses.
% ============================================================
% Tabelle zu §4.2 — Ziel gemäß aktueller Einbindung in chap4.2.tex:
%   tables/chap4/chap4.2table.tex
% Aufruf: % P10-4c, 14.09.2026: Ledger-Herleitung und Nachweisanschlüsse präzisiert.
% B-46-Nachtrag, 13.09.2026: Z3 an die Originalprüfung E43-05
% und die bereits synchronisierten Fassungen von 4.3/Anhang A angeglichen.
% M16 / B-34, 11.09.2026: Z4 trennt die Artefakte der beiden Prüfbeobachtungen.
% M03 / B-17, 10.09.2026: ursprüngliches und erweitertes Nutzungsziel
% unterschieden; keine neue Datierung des Betreuerimpulses.
% ============================================================
% Tabelle zu §4.2 — Ziel gemäß aktueller Einbindung in chap4.2.tex:
%   tables/chap4/chap4.2table.tex
% Aufruf: \input{tables/chap4/chap4.2table}
% Label: t:explorationsverlauf-zyklen
% v03 (01.09., Tabellen-Feedback-Runde — kritisch geprüft): Z1 konkreter
% (Übergaben + Ablaufreihenfolge explizit zu gestalten) · Z2 „als eigene
% Designvariable" · Z4-Designrichtung „an Quellenrepräsentation BINDEN"
% (statt „verlagern") · Z6 ohne MCR-Begriff (fällt erst in 4.5) · Z7
% „deterministische Zuordnungslogik" (VERIFIZIERT am Audit: IngestionApply
% schrieb featureKey-String als Relationsziel — stärkere Variante belegt;
% 4.6-Text + E46-01-Zeile synchronisiert) · Z8 „im Entwicklungsbetrieb"
% (4.7-Grenze schon hier) + „unter Nutzung der bestehenden Governance"
% (= A8-Wortlaut) · Z9 als Stabilisierungs-Übergang · Caption
% „zugehörigen Ankern" (E-IDs+Tags+Kapitelverweise gemischt zulässig).
% v02 (01.09., Kollegen-Review): Zeilen Z0–Z6/Z8/Z9 auf Evidenzgrenzen
% nachgeschärft (keine unbelegten Wertungen/Kausalitäten mehr: Z0 ohne
% „Ergebnisqualität modellabhängig", Z1 als Zielkonkretisierung, Z3
% „unbelegt trotz geschärfter Promptvorgaben", Z4 ohne Zerlegungs-
% Kausalität, Z5 Zirkularität als Eignungsaussage, Z6 Variabilitäts-
% beschreibung statt „inhaltlich scharf", Z8 „regulärer Ablaufzustand",
% Z9 „technische Tests").
% Grundlage: genese-dossier §2 (Z0–Z9), AUDITIERT 31.08. inkl.
% Korrektur-Kopf (A/B/C: zwei ausgeführt, C nie implementiert).
% Systemnamen vermieden; Evidenzanker-Spalte verweist auf die E-IDs
% aus Anhang A und die Artefaktstände (git-Tags) — verbindet Zyklen-
% Überblick und Evidenzsystem.
% ============================================================

\begin{table}[tbp]
\centering
\caption[Verdichteter Explorationsverlauf]{Verdichteter
Explorationsverlauf: Entwicklungszyklen mit prägender Beobachtung
beziehungsweise Zielkonkretisierung, daraus folgender Designrichtung und
den zugehörigen Ankern (E-IDs nach Anhang~\ref{anh:evidenzanker},
markierte Artefaktstände \texttt{v-s}$n$, Kapitelverweise). Eigene
Darstellung.}
\label{t:explorationsverlauf-zyklen}
\scriptsize
\setlength{\tabcolsep}{3pt}
\renewcommand{\arraystretch}{1.2}
\begin{tabularx}{\textwidth}{
    >{\raggedright\arraybackslash}p{0.03\textwidth}
    >{\raggedright\arraybackslash}p{0.16\textwidth}
    >{\raggedright\arraybackslash}X
    >{\raggedright\arraybackslash}p{0.20\textwidth}
    >{\raggedright\arraybackslash}p{0.115\textwidth}
}
\toprule
& \textbf{Fokus} &
\textbf{Prägende Beobachtung / Zielkonkretisierung} &
\textbf{Designrichtung} & \textbf{Anker} \\
\midrule
Z0 & Einzelagent-Durchstich &
Modelltext beschrieb Werkzeugaufrufe ohne registrierte Ausführung &
Ausführung unabhängig von der Modellausgabe beobachtbar machen &
E43-01, E43-02; \texttt{v-s0} \\
\addlinespace[0.3em]
Z1 & linearer Mehr-Agenten-Ablauf &
mit der Aufteilung auf mehrere spezialisierte Rollen mussten auch deren
Informationsübergaben und Ablaufreihenfolge explizit gestaltet werden &
Informationsübergaben und Ablaufsteuerung getrennt von den fachlichen
Rollen betrachten &
\texttt{v-s1} \\
\addlinespace[0.3em]
Z2 & Beobachtbarkeit ausgebaut &
Artefakterzeugung ohne vorausgehenden eigenen Quellzugriff; tatsächlicher
Informationsweg auf Verarbeitungsebene nicht explizit ausgewiesen &
Kontextführung als eigene Designvariable untersuchen &
E43-03; \texttt{v-s2} \\
\addlinespace[0.3em]
Z3 & Strategievergleich (zwei von drei geplanten Varianten) &
nach Promptschärfung zunächst Schreiben ohne Leseaufruf,
im nächsten Lauf Lesen ohne Artefakterzeugung; zwei Varianten
der Kontextführung wurden explorativ verglichen &
Informationsflüsse explizit modellieren &
E43-04, E43-05 \\
\addlinespace[0.3em]
Z4 & nachgelagerte Prüfung als Messinstrument &
freie Prüfung variierte und übersah ausgewählte Inhaltslücken;
festgelegte Prüfeinheiten und Themen ermöglichten gezielte
Einzelprüfungen; die Vollständigkeitsfrage blieb offen &
Vollständigkeitsprüfung an eine explizite Quellenrepräsentation binden &
E44-01, E44-04; \texttt{v-s3} \\
\addlinespace[0.3em]
Z5 & quellgebundene Zwischenrepräsentation &
keine explizite Zuordnung Artefakt--Quellstelle; eine allein
modellgenerierte, nicht menschlich bestätigte Vollständigkeitsreferenz
war als unabhängiger
Bewertungsmaßstab ungeeignet &
Evidenz vor Erzeugung; geteilte Referenzbildung mit menschlicher
Bestätigung &
E44-02, E44-03, E44-05; \texttt{v-s4} \\
\addlinespace[0.3em]
Z6 & getrennte Prüf- und Erzeugungsverantwortung &
freie Prüfung variierte zwischen Wiederholungen in Zahl, Kategorien und
Granularität der Befunde &
Erzeugung, Prüfung und Reparatur als getrennte
Workflow-Verantwortlichkeiten &
E45-02, E45-03; \texttt{v-s5}--\texttt{s6} \\
\addlinespace[0.3em]
Z7 & laufübergreifender Zustand, Kontrolle und Fortschreibung &
Konkretisierung des erweiterten Nutzungsziels (kein Fehlschlag-Pivot);
später: fehlerhafte
Relationen aus einer deterministisch implementierten Zuordnungsregel &
Zustand mit Identitäten, Relationen, Historie; gemeinsamer Mutationspfad
mit Invariantenprüfung; Einordnung neuer Information gegen den Bestand &
E46-01, E47-02; \texttt{v-s7}--\texttt{s8} \\
\addlinespace[0.3em]
Z8 & Außenkopplung, durable Ausführung und Bedienschicht &
menschliche Entscheidungspunkte machen planmäßiges Anhalten zu einem
regulären Ablaufzustand; direkte Einzelbedienung wurde im
Entwicklungsbetrieb mit wachsender Systembreite aufwendiger &
Projektionen aus dem Zustand mit kontrolliertem Rückweg;
Sicherungspunkte und Wiederaufnahme; agentische Bedienebene unter
Nutzung der bestehenden Governance &
E47-01, E48-01/-02; \texttt{v-s9} \\
\addlinespace[0.3em]
Z9 & Stabilisierung des Untersuchungsartefakts &
End-to-End-Abnahmen, Integritätsprüfungen und technische Tests &
stabilisierten Artefaktstand für die Evaluation herstellen &
Kapitel~\ref{chap:maf-realisierung}/\ref{chap:evaluation} \\
\bottomrule
\end{tabularx}
\end{table}

% Label: t:explorationsverlauf-zyklen
% v03 (01.09., Tabellen-Feedback-Runde — kritisch geprüft): Z1 konkreter
% (Übergaben + Ablaufreihenfolge explizit zu gestalten) · Z2 „als eigene
% Designvariable" · Z4-Designrichtung „an Quellenrepräsentation BINDEN"
% (statt „verlagern") · Z6 ohne MCR-Begriff (fällt erst in 4.5) · Z7
% „deterministische Zuordnungslogik" (VERIFIZIERT am Audit: IngestionApply
% schrieb featureKey-String als Relationsziel — stärkere Variante belegt;
% 4.6-Text + E46-01-Zeile synchronisiert) · Z8 „im Entwicklungsbetrieb"
% (4.7-Grenze schon hier) + „unter Nutzung der bestehenden Governance"
% (= A8-Wortlaut) · Z9 als Stabilisierungs-Übergang · Caption
% „zugehörigen Ankern" (E-IDs+Tags+Kapitelverweise gemischt zulässig).
% v02 (01.09., Kollegen-Review): Zeilen Z0–Z6/Z8/Z9 auf Evidenzgrenzen
% nachgeschärft (keine unbelegten Wertungen/Kausalitäten mehr: Z0 ohne
% „Ergebnisqualität modellabhängig", Z1 als Zielkonkretisierung, Z3
% „unbelegt trotz geschärfter Promptvorgaben", Z4 ohne Zerlegungs-
% Kausalität, Z5 Zirkularität als Eignungsaussage, Z6 Variabilitäts-
% beschreibung statt „inhaltlich scharf", Z8 „regulärer Ablaufzustand",
% Z9 „technische Tests").
% Grundlage: genese-dossier §2 (Z0–Z9), AUDITIERT 31.08. inkl.
% Korrektur-Kopf (A/B/C: zwei ausgeführt, C nie implementiert).
% Systemnamen vermieden; Evidenzanker-Spalte verweist auf die E-IDs
% aus Anhang A und die Artefaktstände (git-Tags) — verbindet Zyklen-
% Überblick und Evidenzsystem.
% ============================================================

\begin{table}[tbp]
\centering
\caption[Verdichteter Explorationsverlauf]{Verdichteter
Explorationsverlauf: Entwicklungszyklen mit prägender Beobachtung
beziehungsweise Zielkonkretisierung, daraus folgender Designrichtung und
den zugehörigen Ankern (E-IDs nach Anhang~\ref{anh:evidenzanker},
markierte Artefaktstände \texttt{v-s}$n$, Kapitelverweise). Eigene
Darstellung.}
\label{t:explorationsverlauf-zyklen}
\scriptsize
\setlength{\tabcolsep}{3pt}
\renewcommand{\arraystretch}{1.2}
\begin{tabularx}{\textwidth}{
    >{\raggedright\arraybackslash}p{0.03\textwidth}
    >{\raggedright\arraybackslash}p{0.16\textwidth}
    >{\raggedright\arraybackslash}X
    >{\raggedright\arraybackslash}p{0.20\textwidth}
    >{\raggedright\arraybackslash}p{0.115\textwidth}
}
\toprule
& \textbf{Fokus} &
\textbf{Prägende Beobachtung / Zielkonkretisierung} &
\textbf{Designrichtung} & \textbf{Anker} \\
\midrule
Z0 & Einzelagent-Durchstich &
Modelltext beschrieb Werkzeugaufrufe ohne registrierte Ausführung &
Ausführung unabhängig von der Modellausgabe beobachtbar machen &
E43-01, E43-02; \texttt{v-s0} \\
\addlinespace[0.3em]
Z1 & linearer Mehr-Agenten-Ablauf &
mit der Aufteilung auf mehrere spezialisierte Rollen mussten auch deren
Informationsübergaben und Ablaufreihenfolge explizit gestaltet werden &
Informationsübergaben und Ablaufsteuerung getrennt von den fachlichen
Rollen betrachten &
\texttt{v-s1} \\
\addlinespace[0.3em]
Z2 & Beobachtbarkeit ausgebaut &
Artefakterzeugung ohne vorausgehenden eigenen Quellzugriff; tatsächlicher
Informationsweg auf Verarbeitungsebene nicht explizit ausgewiesen &
Kontextführung als eigene Designvariable untersuchen &
E43-03; \texttt{v-s2} \\
\addlinespace[0.3em]
Z3 & Strategievergleich (zwei von drei geplanten Varianten) &
nach Promptschärfung zunächst Schreiben ohne Leseaufruf,
im nächsten Lauf Lesen ohne Artefakterzeugung; zwei Varianten
der Kontextführung wurden explorativ verglichen &
Informationsflüsse explizit modellieren &
E43-04, E43-05 \\
\addlinespace[0.3em]
Z4 & nachgelagerte Prüfung als Messinstrument &
freie Prüfung variierte und übersah ausgewählte Inhaltslücken;
festgelegte Prüfeinheiten und Themen ermöglichten gezielte
Einzelprüfungen; die Vollständigkeitsfrage blieb offen &
Vollständigkeitsprüfung an eine explizite Quellenrepräsentation binden &
E44-01, E44-04; \texttt{v-s3} \\
\addlinespace[0.3em]
Z5 & quellgebundene Zwischenrepräsentation &
keine explizite Zuordnung Artefakt--Quellstelle; eine allein
modellgenerierte, nicht menschlich bestätigte Vollständigkeitsreferenz
war als unabhängiger
Bewertungsmaßstab ungeeignet &
Evidenz vor Erzeugung; geteilte Referenzbildung mit menschlicher
Bestätigung &
E44-02, E44-03, E44-05; \texttt{v-s4} \\
\addlinespace[0.3em]
Z6 & getrennte Prüf- und Erzeugungsverantwortung &
freie Prüfung variierte zwischen Wiederholungen in Zahl, Kategorien und
Granularität der Befunde &
Erzeugung, Prüfung und Reparatur als getrennte
Workflow-Verantwortlichkeiten &
E45-02, E45-03; \texttt{v-s5}--\texttt{s6} \\
\addlinespace[0.3em]
Z7 & laufübergreifender Zustand, Kontrolle und Fortschreibung &
Konkretisierung des erweiterten Nutzungsziels (kein Fehlschlag-Pivot);
später: fehlerhafte
Relationen aus einer deterministisch implementierten Zuordnungsregel &
Zustand mit Identitäten, Relationen, Historie; gemeinsamer Mutationspfad
mit Invariantenprüfung; Einordnung neuer Information gegen den Bestand &
E46-01, E47-02; \texttt{v-s7}--\texttt{s8} \\
\addlinespace[0.3em]
Z8 & Außenkopplung, durable Ausführung und Bedienschicht &
menschliche Entscheidungspunkte machen planmäßiges Anhalten zu einem
regulären Ablaufzustand; direkte Einzelbedienung wurde im
Entwicklungsbetrieb mit wachsender Systembreite aufwendiger &
Projektionen aus dem Zustand mit kontrolliertem Rückweg;
Sicherungspunkte und Wiederaufnahme; agentische Bedienebene unter
Nutzung der bestehenden Governance &
E47-01, E48-01/-02; \texttt{v-s9} \\
\addlinespace[0.3em]
Z9 & Stabilisierung des Untersuchungsartefakts &
End-to-End-Abnahmen, Integritätsprüfungen und technische Tests &
stabilisierten Artefaktstand für die Evaluation herstellen &
Kapitel~\ref{chap:maf-realisierung}/\ref{chap:evaluation} \\
\bottomrule
\end{tabularx}
\end{table}

% Label: t:explorationsverlauf-zyklen
% v03 (01.09., Tabellen-Feedback-Runde — kritisch geprüft): Z1 konkreter
% (Übergaben + Ablaufreihenfolge explizit zu gestalten) · Z2 „als eigene
% Designvariable" · Z4-Designrichtung „an Quellenrepräsentation BINDEN"
% (statt „verlagern") · Z6 ohne MCR-Begriff (fällt erst in 4.5) · Z7
% „deterministische Zuordnungslogik" (VERIFIZIERT am Audit: IngestionApply
% schrieb featureKey-String als Relationsziel — stärkere Variante belegt;
% 4.6-Text + E46-01-Zeile synchronisiert) · Z8 „im Entwicklungsbetrieb"
% (4.7-Grenze schon hier) + „unter Nutzung der bestehenden Governance"
% (= A8-Wortlaut) · Z9 als Stabilisierungs-Übergang · Caption
% „zugehörigen Ankern" (E-IDs+Tags+Kapitelverweise gemischt zulässig).
% v02 (01.09., Kollegen-Review): Zeilen Z0–Z6/Z8/Z9 auf Evidenzgrenzen
% nachgeschärft (keine unbelegten Wertungen/Kausalitäten mehr: Z0 ohne
% „Ergebnisqualität modellabhängig", Z1 als Zielkonkretisierung, Z3
% „unbelegt trotz geschärfter Promptvorgaben", Z4 ohne Zerlegungs-
% Kausalität, Z5 Zirkularität als Eignungsaussage, Z6 Variabilitäts-
% beschreibung statt „inhaltlich scharf", Z8 „regulärer Ablaufzustand",
% Z9 „technische Tests").
% Grundlage: genese-dossier §2 (Z0–Z9), AUDITIERT 31.08. inkl.
% Korrektur-Kopf (A/B/C: zwei ausgeführt, C nie implementiert).
% Systemnamen vermieden; Evidenzanker-Spalte verweist auf die E-IDs
% aus Anhang A und die Artefaktstände (git-Tags) — verbindet Zyklen-
% Überblick und Evidenzsystem.
% ============================================================

\begin{table}[tbp]
\centering
\caption[Verdichteter Explorationsverlauf]{Verdichteter
Explorationsverlauf: Entwicklungszyklen mit prägender Beobachtung
beziehungsweise Zielkonkretisierung, daraus folgender Designrichtung und
den zugehörigen Ankern (E-IDs nach Anhang~\ref{anh:evidenzanker},
markierte Artefaktstände \texttt{v-s}$n$, Kapitelverweise). Eigene
Darstellung.}
\label{t:explorationsverlauf-zyklen}
\scriptsize
\setlength{\tabcolsep}{3pt}
\renewcommand{\arraystretch}{1.2}
\begin{tabularx}{\textwidth}{
    >{\raggedright\arraybackslash}p{0.03\textwidth}
    >{\raggedright\arraybackslash}p{0.16\textwidth}
    >{\raggedright\arraybackslash}X
    >{\raggedright\arraybackslash}p{0.20\textwidth}
    >{\raggedright\arraybackslash}p{0.115\textwidth}
}
\toprule
& \textbf{Fokus} &
\textbf{Prägende Beobachtung / Zielkonkretisierung} &
\textbf{Designrichtung} & \textbf{Anker} \\
\midrule
Z0 & Einzelagent-Durchstich &
Modelltext beschrieb Werkzeugaufrufe ohne registrierte Ausführung &
Ausführung unabhängig von der Modellausgabe beobachtbar machen &
E43-01, E43-02; \texttt{v-s0} \\
\addlinespace[0.3em]
Z1 & linearer Mehr-Agenten-Ablauf &
mit der Aufteilung auf mehrere spezialisierte Rollen mussten auch deren
Informationsübergaben und Ablaufreihenfolge explizit gestaltet werden &
Informationsübergaben und Ablaufsteuerung getrennt von den fachlichen
Rollen betrachten &
\texttt{v-s1} \\
\addlinespace[0.3em]
Z2 & Beobachtbarkeit ausgebaut &
Artefakterzeugung ohne vorausgehenden eigenen Quellzugriff; tatsächlicher
Informationsweg auf Verarbeitungsebene nicht explizit ausgewiesen &
Kontextführung als eigene Designvariable untersuchen &
E43-03; \texttt{v-s2} \\
\addlinespace[0.3em]
Z3 & Strategievergleich (zwei von drei geplanten Varianten) &
nach Promptschärfung zunächst Schreiben ohne Leseaufruf,
im nächsten Lauf Lesen ohne Artefakterzeugung; zwei Varianten
der Kontextführung wurden explorativ verglichen &
Informationsflüsse explizit modellieren &
E43-04, E43-05 \\
\addlinespace[0.3em]
Z4 & nachgelagerte Prüfung als Messinstrument &
freie Prüfung variierte und übersah ausgewählte Inhaltslücken;
festgelegte Prüfeinheiten und Themen ermöglichten gezielte
Einzelprüfungen; die Vollständigkeitsfrage blieb offen &
Vollständigkeitsprüfung an eine explizite Quellenrepräsentation binden &
E44-01, E44-04; \texttt{v-s3} \\
\addlinespace[0.3em]
Z5 & quellgebundene Zwischenrepräsentation &
keine explizite Zuordnung Artefakt--Quellstelle; eine allein
modellgenerierte, nicht menschlich bestätigte Vollständigkeitsreferenz
war als unabhängiger
Bewertungsmaßstab ungeeignet &
Evidenz vor Erzeugung; geteilte Referenzbildung mit menschlicher
Bestätigung &
E44-02, E44-03, E44-05; \texttt{v-s4} \\
\addlinespace[0.3em]
Z6 & getrennte Prüf- und Erzeugungsverantwortung &
freie Prüfung variierte zwischen Wiederholungen in Zahl, Kategorien und
Granularität der Befunde &
Erzeugung, Prüfung und Reparatur als getrennte
Workflow-Verantwortlichkeiten &
E45-02, E45-03; \texttt{v-s5}--\texttt{s6} \\
\addlinespace[0.3em]
Z7 & laufübergreifender Zustand, Kontrolle und Fortschreibung &
Konkretisierung des erweiterten Nutzungsziels (kein Fehlschlag-Pivot);
später: fehlerhafte
Relationen aus einer deterministisch implementierten Zuordnungsregel &
Zustand mit Identitäten, Relationen, Historie; gemeinsamer Mutationspfad
mit Invariantenprüfung; Einordnung neuer Information gegen den Bestand &
E46-01, E47-02; \texttt{v-s7}--\texttt{s8} \\
\addlinespace[0.3em]
Z8 & Außenkopplung, durable Ausführung und Bedienschicht &
menschliche Entscheidungspunkte machen planmäßiges Anhalten zu einem
regulären Ablaufzustand; direkte Einzelbedienung wurde im
Entwicklungsbetrieb mit wachsender Systembreite aufwendiger &
Projektionen aus dem Zustand mit kontrolliertem Rückweg;
Sicherungspunkte und Wiederaufnahme; agentische Bedienebene unter
Nutzung der bestehenden Governance &
E47-01, E48-01/-02; \texttt{v-s9} \\
\addlinespace[0.3em]
Z9 & Stabilisierung des Untersuchungsartefakts &
End-to-End-Abnahmen, Integritätsprüfungen und technische Tests &
stabilisierten Artefaktstand für die Evaluation herstellen &
Kapitel~\ref{chap:maf-realisierung}/\ref{chap:evaluation} \\
\bottomrule
\end{tabularx}
\end{table}

% Label: t:explorationsverlauf-zyklen
% v03 (01.09., Tabellen-Feedback-Runde — kritisch geprüft): Z1 konkreter
% (Übergaben + Ablaufreihenfolge explizit zu gestalten) · Z2 „als eigene
% Designvariable" · Z4-Designrichtung „an Quellenrepräsentation BINDEN"
% (statt „verlagern") · Z6 ohne MCR-Begriff (fällt erst in 4.5) · Z7
% „deterministische Zuordnungslogik" (VERIFIZIERT am Audit: IngestionApply
% schrieb featureKey-String als Relationsziel — stärkere Variante belegt;
% 4.6-Text + E46-01-Zeile synchronisiert) · Z8 „im Entwicklungsbetrieb"
% (4.7-Grenze schon hier) + „unter Nutzung der bestehenden Governance"
% (= A8-Wortlaut) · Z9 als Stabilisierungs-Übergang · Caption
% „zugehörigen Ankern" (E-IDs+Tags+Kapitelverweise gemischt zulässig).
% v02 (01.09., Kollegen-Review): Zeilen Z0–Z6/Z8/Z9 auf Evidenzgrenzen
% nachgeschärft (keine unbelegten Wertungen/Kausalitäten mehr: Z0 ohne
% „Ergebnisqualität modellabhängig", Z1 als Zielkonkretisierung, Z3
% „unbelegt trotz geschärfter Promptvorgaben", Z4 ohne Zerlegungs-
% Kausalität, Z5 Zirkularität als Eignungsaussage, Z6 Variabilitäts-
% beschreibung statt „inhaltlich scharf", Z8 „regulärer Ablaufzustand",
% Z9 „technische Tests").
% Grundlage: genese-dossier §2 (Z0–Z9), AUDITIERT 31.08. inkl.
% Korrektur-Kopf (A/B/C: zwei ausgeführt, C nie implementiert).
% Systemnamen vermieden; Evidenzanker-Spalte verweist auf die E-IDs
% aus Anhang A und die Artefaktstände (git-Tags) — verbindet Zyklen-
% Überblick und Evidenzsystem.
% ============================================================

\begin{table}[tbp]
\centering
\caption[Verdichteter Explorationsverlauf]{Verdichteter
Explorationsverlauf: Entwicklungszyklen mit prägender Beobachtung
beziehungsweise Zielkonkretisierung, daraus folgender Designrichtung und
den zugehörigen Ankern (E-IDs nach Anhang~\ref{anh:evidenzanker},
markierte Artefaktstände \texttt{v-s}$n$, Kapitelverweise). Eigene
Darstellung.}
\label{t:explorationsverlauf-zyklen}
\scriptsize
\setlength{\tabcolsep}{3pt}
\renewcommand{\arraystretch}{1.2}
\begin{tabularx}{\textwidth}{
    >{\raggedright\arraybackslash}p{0.03\textwidth}
    >{\raggedright\arraybackslash}p{0.16\textwidth}
    >{\raggedright\arraybackslash}X
    >{\raggedright\arraybackslash}p{0.20\textwidth}
    >{\raggedright\arraybackslash}p{0.115\textwidth}
}
\toprule
& \textbf{Fokus} &
\textbf{Prägende Beobachtung / Zielkonkretisierung} &
\textbf{Designrichtung} & \textbf{Anker} \\
\midrule
Z0 & Einzelagent-Durchstich &
Modelltext beschrieb Werkzeugaufrufe ohne registrierte Ausführung &
Ausführung unabhängig von der Modellausgabe beobachtbar machen &
E43-01, E43-02; \texttt{v-s0} \\
\addlinespace[0.3em]
Z1 & linearer Mehr-Agenten-Ablauf &
mit der Aufteilung auf mehrere spezialisierte Rollen mussten auch deren
Informationsübergaben und Ablaufreihenfolge explizit gestaltet werden &
Informationsübergaben und Ablaufsteuerung getrennt von den fachlichen
Rollen betrachten &
\texttt{v-s1} \\
\addlinespace[0.3em]
Z2 & Beobachtbarkeit ausgebaut &
Artefakterzeugung ohne vorausgehenden eigenen Quellzugriff; tatsächlicher
Informationsweg auf Verarbeitungsebene nicht explizit ausgewiesen &
Kontextführung als eigene Designvariable untersuchen &
E43-03; \texttt{v-s2} \\
\addlinespace[0.3em]
Z3 & Strategievergleich (zwei von drei geplanten Varianten) &
nach Promptschärfung zunächst Schreiben ohne Leseaufruf,
im nächsten Lauf Lesen ohne Artefakterzeugung; zwei Varianten
der Kontextführung wurden explorativ verglichen &
Informationsflüsse explizit modellieren &
E43-04, E43-05 \\
\addlinespace[0.3em]
Z4 & nachgelagerte Prüfung als Messinstrument &
freie Prüfung variierte und übersah ausgewählte Inhaltslücken;
festgelegte Prüfeinheiten und Themen ermöglichten gezielte
Einzelprüfungen; die Vollständigkeitsfrage blieb offen &
Vollständigkeitsprüfung an eine explizite Quellenrepräsentation binden &
E44-01, E44-04; \texttt{v-s3} \\
\addlinespace[0.3em]
Z5 & quellgebundene Zwischenrepräsentation &
keine explizite Zuordnung Artefakt--Quellstelle; eine allein
modellgenerierte, nicht menschlich bestätigte Vollständigkeitsreferenz
war als unabhängiger
Bewertungsmaßstab ungeeignet &
Evidenz vor Erzeugung; geteilte Referenzbildung mit menschlicher
Bestätigung &
E44-02, E44-03, E44-05; \texttt{v-s4} \\
\addlinespace[0.3em]
Z6 & getrennte Prüf- und Erzeugungsverantwortung &
freie Prüfung variierte zwischen Wiederholungen in Zahl, Kategorien und
Granularität der Befunde &
Erzeugung, Prüfung und Reparatur als getrennte
Workflow-Verantwortlichkeiten &
E45-02, E45-03; \texttt{v-s5}--\texttt{s6} \\
\addlinespace[0.3em]
Z7 & laufübergreifender Zustand, Kontrolle und Fortschreibung &
Konkretisierung des erweiterten Nutzungsziels (kein Fehlschlag-Pivot);
später: fehlerhafte
Relationen aus einer deterministisch implementierten Zuordnungsregel &
Zustand mit Identitäten, Relationen, Historie; gemeinsamer Mutationspfad
mit Invariantenprüfung; Einordnung neuer Information gegen den Bestand &
E46-01, E47-02; \texttt{v-s7}--\texttt{s8} \\
\addlinespace[0.3em]
Z8 & Außenkopplung, durable Ausführung und Bedienschicht &
menschliche Entscheidungspunkte machen planmäßiges Anhalten zu einem
regulären Ablaufzustand; direkte Einzelbedienung wurde im
Entwicklungsbetrieb mit wachsender Systembreite aufwendiger &
Projektionen aus dem Zustand mit kontrolliertem Rückweg;
Sicherungspunkte und Wiederaufnahme; agentische Bedienebene unter
Nutzung der bestehenden Governance &
E47-01, E48-01/-02; \texttt{v-s9} \\
\addlinespace[0.3em]
Z9 & Stabilisierung des Untersuchungsartefakts &
End-to-End-Abnahmen, Integritätsprüfungen und technische Tests &
stabilisierten Artefaktstand für die Evaluation herstellen &
Kapitel~\ref{chap:maf-realisierung}/\ref{chap:evaluation} \\
\bottomrule
\end{tabularx}
\end{table}
